\title{FlashAttention-3: Fast and Accurate Attention with Asynchrony and Low-precision}

\begin{abstract}
The attention mechanism, central to the Transformer architecture, often constitutes the primary computational bottleneck in large language models and applications requiring extended contexts. \textsc{FlashAttention} introduced a strategy for accelerating attention calculations on GPUs by reducing memory traffic, yet this method does not fully leverage advancements in the latest hardware generations; notably, \textsc{FlashAttention-2} only attains 35\% hardware utilization on the H100 GPU. In this work, we present three key innovations to further optimize attention operations on Hopper GPUs: we utilize the asynchronous nature of Tensor Cores and TMA to (1) enable concurrent computation and data transfer via warp-specialization, and (2) intertwine block-wise matrix multiplication and softmax computations, as well as (3) employ block quantization and incoherent execution strategies facilitated by the hardware's FP8 low-precision support. Our proposed approach, \textsc{FlashAttention-3}, provides a performance increase of 1.5-2.0$\times$ on H100 GPUs, achieving up to 740 TFLOPs/s with FP16 (reaching 75\% utilization), and approaching 1.2 PFLOPs/s with FP8 precision. We further show that FP8 \textsc{FlashAttention-3} incurs 2.6$\times$ less numerical error compared to a baseline FP8 attention implementation.
\end{abstract}

\section{Introduction}
\label{sec:intro}

Within the Transformer architecture~\citep{vaswani2017attention}, the most significant computational limitation stems from the attention mechanism, as calculating the self-attention between queries and keys requires computation that increases quadratically with sequence length. Enhancing the scalability of attention to accommodate longer contexts can facilitate advanced modeling and reasoning tasks involving extended documents~\citep{guo2021longt5,shaham2022scrolls,peng2023yarn}, extensive code repositories~\citep{roziere2023code, li2023starcoder}, as well as enable processing of novel modalities such as high-resolution imagery~\citep{chen2022scaling}, audio~\citep{gulati2020conformer}, and video~\citep{ho2022video}. Furthermore, applications like maintaining detailed user interaction histories~\citep{sun2019bert4rec} and managing agent activities with prolonged time horizons~\citep{yao2022react} become feasible. Consequently, there is considerable focus on expediting attention mechanisms in the context of long sequences, employing strategies such as approximate attention methods~\citep{katharopoulos2020transformers,choromanski2020rethinking, tay2020efficient}, software-level optimizations~(\citep{rabe2021self,dao2022flashattention,kwon2023efficient}), and exploring alternative architectural designs~\citep{peng2023rwkv,sun2023retentive,gu2023mamba}.

In this study, we extend the efforts of \citet{dao2022flashattention}, who pioneered the design of exact-attention algorithms by leveraging insights into GPU hardware architecture and execution models at an early design stage. Dao et al. \citep{dao2022flashattention} presented \textsc{FlashAttention}, introducing an innovative tiling approach for parallelizing attention computations that avoids the need for intermediate accesses to slow global memory by merging all attention-related operations within a single GPU kernel. Building upon this, \citet{dao2023flashattention2} reformulated the algorithm as \textsc{FlashAttention-2}, enabling parallelization across the sequence length and executing the inner loop of the forward pass over segmented blocks of the key and value matrices, thereby enhancing workload distribution and occupancy on the GPU. Nevertheless, we notice that \textsc{FlashAttention-2} attains significantly lower hardware utilization on newer GPUs compared to state-of-the-art matrix-multiplication (GEMM) kernels, with rates such as 35\% versus 80-90\% on the Hopper H100 GPU. This suboptimal performance can be partly traced to implementation differences, including the lack of Hopper-specific instruction utilization, instead relying on instructions optimized for earlier Ampere Tensor Cores. Recent advances, such as ThunkerKitten~\citep{spector2024thunder} and cuDNN 9~\citep{cudnn9}, have demonstrated that incorporating Hopper-targeted instructions and tile-oriented abstractions not only accelerates attention computations but also streamlines code development.

At a more basic level, the design of \textsc{FlashAttention-2}'s algorithm is rooted in a simplified, synchronous execution model and does not incorporate explicit mechanisms for leveraging asynchrony or low-precision arithmetic. Asynchrony emerges due to specialized hardware architectures that accelerate key components of machine learning workloads—namely, dedicated hardware units such as Tensor Cores for matrix multiplication and Tensor Memory Accelerator (TMA) for data movement, both of which function independently from the general-purpose CUDA cores that handle logic and arithmetic operations. The use of low-precision formats—FP8 on Hopper and FP4 on Blackwell, following the earlier adoption of FP16 (Pascal, 2017) and BF16 (Ampere, 2020)—continues to demonstrate increased throughput (2$\times$ or 4$\times$) without expanding power consumption or silicon area. In \cref{subsec:hardware}, we summarize Hopper's offerings in these respects. 
The main technical obstacle involves reengineering \textsc{FlashAttention-2} to exploit such hardware functionalities: supporting asynchrony entails carefully orchestrating overlaps between matrix multiplication and softmax operations—even with existing dependencies—while adapting to low-precision operations requires strategies to contain quantization errors, which is particularly important given the occurrence of outlier features in LLMs~\citep{dettmers2208llm, sun2024massive}.

With this objective, we introduce \textsc{FlashAttention-3}, which brings together three novel approaches aimed at further optimizing performance on the latest GPU architectures.\footnote{Our findings are presented primarily in relation to NVIDIA's Hopper architecture, though the proposed algorithm is applicable to any GPU platform featuring strong asynchronous execution and support for low-precision operations.}

\begin{enumerate}

\item \textbf{Producer-Consumer asynchrony:} We propose a warp-specialized software pipelining approach that leverages the ability to decouple the execution of data transfers and Tensor Core computations. By allocating data producers and consumers to distinct warps, our method enhances the algorithm's potential to mask both memory access and instruction issue delays.
\item \textbf{Hiding softmax under asynchronous block-wise GEMMs:} To mask the relatively low-throughput non-GEMM computations required in softmax—such as floating-point multiply-add and exponential operations—we schedule them concurrently with asynchronous WGMMA GEMM instructions. This entails a re-architecting of the \textsc{FlashAttention-2} algorithm to break select sequential dependencies between softmax and the GEMM operations. For instance, in our two-stage algorithm, while softmax is being computed on a particular block of the scores matrix, WGMMA performs an asynchronous prefetch and execution for the next block.
\item \textbf{Hardware-accelerated low-precision GEMM:} We modify the forward pass algorithm to target FP8 Tensor Cores for GEMM computations, which nearly doubles the observed throughput in TFLOPs/s. Accommodating the differing memory layout conventions required by WGMMA for blocks of FP32 accumulators and FP8 input matrices involves resolving layout mismatches. To address potential accuracy degradation from employing FP8 precision, we integrate block quantization and incoherent processing strategies.
\end{enumerate}

For empirical evaluation, we conduct extensive benchmarking of \textsc{FlashAttention-3} on the H100 SXM5 GPU across various parameter settings. The results demonstrate that (1) in FP16 precision, \textsc{FlashAttention-3} attains a 1.5–2.0$\times$ acceleration in the forward pass relative to \textsc{FlashAttention-2}—reaching peak throughput of up to 740 TFLOPs/s—and a 1.5–1.75$\times$ speedup in the backward pass, (2) in FP8 precision, the throughput approaches 1.2 PFLOPs/s, and (3) for longer sequence lengths, FP16 surpasses, and FP8 remains competitive with, the top-performing attention implementation from NVIDIA's cuDNN library.\footnote{Specifically, for head dimension 64, FP8 \textsc{FlashAttention-3} leads, while for head dimensions 128 and 256, FP8 performance matches cuDNN without causal masking and slightly trails it under causal masking.} Our assessment further confirms that FP16 \textsc{FlashAttention-3} preserves numerical accuracy on par with \textsc{FlashAttention-2}, and exceeds the baseline attention implementation in accuracy, owing to all intermediate computations (e.g., softmax scaling) being maintained in FP32. Additionally, FP8 \textsc{FlashAttention-3}, utilizing block quantization and incoherent processing, delivers 2.6$\times$ improved accuracy over the baseline per-tensor quantization approach in scenarios with outlier features.

We have released \textsc{FlashAttention-3} as open-source software under a permissive license\footnote{\textsc{FlashAttention-3}
  is available at \texttt{https://github.com/Dao-AILab/flash-attention}}, and intend to incorporate it into the PyTorch and Hugging Face ecosystems, aiming to maximize its accessibility and usefulness for researchers and practitioners.

\section{Background: Multi-Head Attention and GPU Characteristics}
\label{sec:background}

\subsection{Multi-Head Attention}
\label{subsec:multi_head_attn}

Consider the input sequences for queries, keys, and values, denoted as $\mathbf{Q}, \mathbf{K}, \mathbf{V} \in \mathbb{R}^{N \times d}$, each pertaining to one attention head, where $N$ represents the length of the input sequence and $d$ denotes the dimensionality of the head. The corresponding attention output $\mathbf{O}$ can then be determined using the following procedure:

\begin{equation*}
  \mathbf{S} = \alpha \mathbf{Q} \mathbf{K}^\top \in \mathbb{R}^{N \times N}, \quad \mathbf{P} = \mathrm{softmax}(\mathbf{S}) \in \mathbb{R}^{N \times N}, \quad \mathbf{O} = \mathbf{P}\mathbf{V} \in \mathbb{R}^{N \times d},
\end{equation*}

Here, $\mathrm{softmax}$ operates across each row, with the scaling parameter commonly chosen as $\alpha = 1/\sqrt{d}$. To enhance numerical stability during exponentiation, it is standard practice to subtract $\mathrm{rowmax}(\mathbf{S})$ from $\mathbf{S}$. In the context of multi-head attention (MHA), separate projection matrices for queries, keys, and values are assigned to each attention head. These computations are performed concurrently across different heads and batches, yielding the complete output tensor.

Consider $\phi$ as a scalar loss function, and denote the gradient operator as $\mathbf{d}(-) = \partial \phi / \partial (-)$. With the output gradient $\mathbf{dO} \in \mathbb{R}^{N \times d}$ provided, we apply the chain rule to derive $\mathbf{dQ}$, $\mathbf{dK}$, and $\mathbf{dV}$ as outlined below:

\begin{align*}
  \mathbf{dV} &= \mathbf{P}^\top \mathbf{dO} \in \mathbb{R}^{N \times d} \\
  \mathbf{dP} &= \mathbf{dO} \mathbf{V}^\top \in \mathbb{R}^{N \times N} \\
  \mathbf{dS} &= \mathrm{dsoftmax} (\mathbf{dP}) \in \mathbb{R}^{N \times N} \\
  \mathbf{dQ} &= \alpha \mathbf{dS} \mathbf{K} \in \mathbb{R}^{N \times d} \\
  \mathbf{dK} &= \alpha \mathbf{dS}^\top \mathbf{Q} \in \mathbb{R}^{N \times d},
\end{align*}

In this context, the relationship $\mathbf{d}s = (\mathrm{diag}(p) - p p^\top)\mathbf{d}p$ holds for $p = \mathrm{softmax}(s)$ when expressed as a function of the vector $s$, and we denote the row-wise application of this operation by $\mathrm{dsoftmax}(\mathbf{dP})$. This calculation is also amenable to parallelization over both the heads and batch dimensions during the backward pass of MHA.

\subsection{GPU hardware characteristics and execution model}
\label{subsec:hardware}

We outline the pertinent components of the GPU execution model that impact \textsc{FlashAttention-3}, concentrating on the NVIDIA Hopper architecture as a specific example of this framework.

\paragraph{Memory hierarchy:} The memory system of the GPU is structured as a hierarchy of storage locations, where higher bandwidth typically corresponds to lower storage capacity (\cref{tab:gpu-hierarchy})\footnote{\citet{luo2024benchmarking} provides shared memory bandwidth as 128 bytes per clock cycle for each SM. We obtain total bandwidth by multiplying this by 132 SMs and a boost clock frequency of 1830 MHz.}. At the top level, global memory (GMEM)—alternatively called HBM—refers to off-chip DRAM that is universally accessible to all streaming multiprocessors (SMs). Accesses to GMEM are automatically cached in the on-chip L2 cache. At the next level, every SM has its own on-chip, highly banked shared memory (SMEM), which is directly managed by programmers. At the lowest level of the hierarchy sits the register file specific to each SM.

\paragraph{Thread hierarchy:} The structure of GPU programming is built on various levels of grouped execution units, collectively referred to as threads. At the most granular level are individual threads, which are assembled into warps, each containing 32 threads. Four sequential warps form a warpgoup. Above this are threadblocks—also known as cooperative thread arrays (CTAs)—which can be further combined into threadblock clusters, a level introduced in Hopper architectures, and the highest level consists of grids.

The two hierarchies exhibit a strong interconnection. Threads belonging to a single CTA are executed concurrently on the same SM, while CTAs grouped within one cluster are jointly scheduled onto the same GPC. All threads inside a CTA can directly access SMEM, in contrast to RMEM, where each thread possesses a maximum of 256 private registers.

\paragraph{Asynchrony and warp-specialization:}

GPUs function as throughput-oriented processors that utilize parallelism and asynchronous operations to mask delays arising from memory accesses and instruction execution. To facilitate asynchronous memory transfers between GMEM and SMEM, Hopper introduces a specialized hardware module known as the Tensor Memory Accelerator (TMA) \cite[\S7.29]{cuda}. Moreover, differentiating itself from older architectures like Ampere, Hopper’s Tensor Core—accessible through the warpgroup-scoped WGMMA instruction \cite[\S9.7.14]{ptx}—operates asynchronously and is capable of directly obtaining its input data from shared memory.

The incorporation of hardware mechanisms for asynchrony enables the use of warp-specialized kernels, wherein the warps within a CTA are assigned distinct functions as either producers or consumers, restricting them to performing solely data transfer or computational tasks, respectively. This division generally enhances the compiler’s effectiveness in constructing optimal instruction schedules \citep{warp-specialization-2011}. Furthermore, Hopper facilitates on-the-fly redistribution of registers among warpgroups using the \verb|setmaxnreg| command \cite[\S9.7.17.1]{ptx}, permitting those warps responsible for executing MMAs to be allocated a greater portion of RMEM compared to those limited to TMA operations (which typically require only one thread).

\paragraph{Low-precision number formats:}
\label{sec:low-precision-gpu}
Recent generations of GPUs incorporate dedicated hardware units specifically designed for efficient low-precision arithmetic. As an illustration, Hopper architecture’s FP8 Tensor Cores, accessed via the WGMMA instruction, can achieve double the throughput per SM relative to computations using FP16 or BF16 precision.

Nonetheless, effective utilization of FP8 WGMMA requires familiarity with the specific memory layout restrictions imposed on its input operands. Consider a GEMM operation that computes $A \times B^{\top}$, where $A$ is an $M\times K$ matrix and $B$ is an $N\times K$ matrix. An operand is designated as \emph{mn-major} when it is stored contiguously along the outer $M$ or $N$ axis, and as \emph{k-major} when contiguity is along the inner $K$ axis. With FP16 WGMMA, both mn-major and k-major representations of operands held in SMEM are valid. In contrast, FP8 WGMMA is limited to supporting only operands in k-major format. Furthermore, in compound scenarios—such as in attention mechanisms—where multiple GEMMs are fused within one kernel, mismatches between the layout of FP32 accumulators and FP8 operands can prevent the use of chained, dependent FP8 WGMMAs.

With regard to attention mechanisms, these layout constraints require specific adjustments to the structure of an FP8 algorithm, as detailed in \cref{sec:algofp8}.

\subsection{Standard Attention and Flash Attention} In line with \citet{dao2022flashattention}, we use the term \textbf{standard attention} to refer to the conventional GPU-based attention computation, where the intermediate matrices $\mathbf{S}$ and $\mathbf{P}$ are explicitly stored in high-bandwidth memory (HBM). \textsc{FlashAttention} introduces a different approach by employing a localized softmax reduction. This technique eliminates the need for costly intermediate memory operations by consolidating the entire attention computation into a single kernel. The localized softmax aligns with lines \ref{code-ws:softmax_start}-\ref{code-ws:softmax_end} of the consumer mainloop described in \cref{alg:flash3_wgmma_ws_only}, in conjunction with the rescaling steps applied to the blocks of $\mathbf{O}$. A straightforward argument demonstrating that this method yields the correct computation of $\mathbf{O}$ is presented in \cite[\S 2.3.1]{dao2023flashattention2}.

\section{FlashAttention-3: Algorithm}
\label{sec:algo}

This section presents an overview of the \textsc{FlashAttention-3} algorithm. Our discussion primarily centers on the forward pass, while details of the backward pass can be found in~\cref{sec:algo_ws_bwd}. We begin by outlining the integration of warp-specialization and a circular SMEM buffer into the core structure of \textsc{FlashAttention-2}. Next, we discuss leveraging WGMMA's asynchronous capabilities to create a two-stage pipeline in which GEMM and softmax operations are overlapped. Lastly, we address the adaptations necessary for FP8, considering both proper layout alignment and enhancing precision through block quantization as well as handling incoherent data processing.

\subsection{Producer-Consumer asynchrony through warp-specialization and pingpong scheduling}
\label{sec:algo_ws}

\paragraph{Warp-specialization}
Similar to \textsc{FlashAttention-2}, the forward computation in \textsc{FlashAttention-3} is highly parallelizable across dimensions such as batch size, number of attention heads, and the length of the query sequence. Accordingly, presenting the algorithm from the perspective of a CTA is sufficient; this approach processes a query tile $\mathbf{Q}_i$ and produces its associated output tile $\mathbf{O}_i$. For clarity, we begin by outlining the warp-specialization approach that utilizes a circular SMEM buffer, while excluding the GEMM-softmax overlap from consideration. Let $d$ denote the head dimension, and $N$ the length of the sequence. By selecting a query block size $B_r$, the query matrix $\mathbf{Q}$ is partitioned into $T_r = \lceil \frac{N}{B_r} \rceil$ blocks $\mathbf{Q}_1, .., \mathbf{Q}_{T_r}$.

\begin{algorithm}[H]
    \caption{\small\label{alg:flash3_wgmma_ws_only}\textsc{FlashAttention-3} forward pass \textbf{excluding} intra-consumer overlapping -- CTA-centric approach}
    \begin{algorithmic}[1]
\REQUIRE Input matrices $\mathbf{Q}_i \in \mathbb{R}^{B_r \times d}$, $\mathbf{K}, \mathbf{V} \in \mathbb{R}^{N \times d}$ located in HBM; key block dimension $B_c$ and corresponding count $T_c = \lceil \frac{N}{B_c} \rceil$.
\STATE Set up pipeline manager responsible for barrier synchronization with a circular SMEM buffer spanning $s$ stages.
\IF {current warpgroup is designated producer}
\STATE Free a preset quantity of registers.
\STATE Load $\mathbf{Q}_i$ from HBM into shared memory.
\STATE On successful load, update status to inform consumers about the availability of $\mathbf{Q}_i$.
\FOR{$0 \le j < T_c$}
    \STATE Block until the $(j\,\%\,s)$ stage of the buffer is vacated by the consumer.
    \STATE Transfer $\mathbf{K}_j, \mathbf{V}_j$ from HBM into the $(j\,\%\,s)$ stage of shared memory buffer.
    \STATE Once loading finishes, update status to notify consumers about $\mathbf{K}_j, \mathbf{V}_j$ readiness.
\ENDFOR
\ELSE \STATE Allocate a preset number of registers according to the total consumer warps.
\STATE Initialize, on-chip, $\mathbf{O}_i = (0) \in \mathbb{R}^{B_r \times d}$ as well as $\ell_i = 0$ and $m_i = -\infty$ within $\mathbb{R}^{B_r}$.
\STATE Hold until $\mathbf{Q}_i$ appears in shared memory.
\FOR{$0 \le j < T_c$}
\STATE Wait for $\mathbf{K}_j$'s transfer into shared memory to complete.
\STATE Execute the SS-GEMM: $\mathbf{S}_i^{(j)} = \mathbf{Q}_i \mathbf{K}_j^T$. Commit, then wait. \label{alg:ws_only_gemm1}
\STATE Save current $m_i$ as $m_i^{\mathrm{old}}$, then update $m_i$ to $\mathrm{max}(m_i^{\mathrm{old}}, \mathrm{rowmax}(\mathbf{S}_i^{(j)}))$. \label{code-ws:softmax_start}
\STATE Calculate $\widetilde{\mathbf{P}}_i^{(j)} = \mathrm{exp}(\mathbf{S}_i^{(j)} - m_i)$ and refresh $\ell_i = \mathrm{exp}(m_i^{\mathrm{old}} - m_i) \ell_i + \mathrm{rowsum}(\widetilde{\mathbf{P}}_i^{(j)})$. \label{code-ws:softmax_end}
\STATE Wait for $\mathbf{V}_j$ to be fully available in shared memory.
\STATE Update $\mathbf{O}_i = \mathrm{diag}(\mathrm{exp}(m_i^{\mathrm{old}} - m_i))^{-1} \mathbf{O}_i + \widetilde{\mathbf{P}}_i^{(j)} \mathbf{V}_j$ through RS-GEMM. Commit and await. \label{alg:ws_only_gemm2}
\STATE Signal the producer by releasing the $(j\,\%\,s)$ stage of the shared buffer.
\ENDFOR
\STATE Compute the normalization: $\mathbf{O}_i = \mathrm{diag}(\ell_i)^{-1} \mathbf{O}_i$; calculate $L_i = m_i + \log(\ell_i)$.
\STATE Store the results, writing $\mathbf{O}_i$ and $L_i$ to HBM in the positions of the $i$th block of $\mathbf{O}$ and $L$.
\ENDIF
\end{algorithmic}
\end{algorithm}

In our implementation of \cref{alg:flash3_wgmma_ws_only} targeting the Hopper architecture, we employ \verb|setmaxnreg| to manage register allocation and deallocation. The transfer of $\mathbf{Q}_i$ and $\{ \mathbf{K}_j, \mathbf{V}_j \}_{0 \leq j < T_c}$ is performed via TMA loads. The GEMMs within the consumer mainloop utilize WGMMA, where the prefixes SS or RS specify whether the first input operand resides in shared memory or the register file, respectively. It is important to recognize that, according to the execution structure of \cref{alg:flash3_wgmma_ws_only}, TMA load instructions proceed asynchronously; issuing these loads does not require prior loads to complete. Furthermore, the producer mainloop will not introduce any waits during the initial $s$ iterations while the buffer is being populated.

\paragraph{Pingpong scheduling} The asynchronous behavior of WGMMA and TMA, when combined with warp-specialization, enables concurrent execution such that while one warpgroup performs softmax computations, another can execute GEMM tasks. This strategy is particularly beneficial given the stark disparity in throughput between matmul and non-matmul instructions on contemporary accelerators. For instance, the H100 SXM5 GPU achieves 989 TFLOPS in FP16 matmul, but merely 3.9 TFLOPS for special operations such as exponentials\footnote{According to the CUDA programming guide, each streaming multiprocessor (SM) can execute 16 special function operations per clock cycle. This yields 3.9 TFLOPS when calculating $16 \times 132$ SMs at 1830 MHz.}. In the context of an attention forward pass in FP16 with a head size of 128, the workload comprises 512 times more matmul FLOPS relative to exponential operations, despite exponentials offering 256 times less throughput; thus, exponentials may occupy up to 50\% of the compute cycles needed for matmuls. This imbalance becomes even more pronounced in FP8, where matmul performance doubles but exponential throughput remains unchanged.

Given that the exponential operation is handled by a dedicated hardware component—the multi-function unit—it is preferable to coordinate the exponential computation to occur concurrently with the matrix multiplication executed by the Tensor Cores. To achieve this concurrency, we introduce synchronization barriers (\texttt{bar.sync} instructions), thereby enforcing an ordering in which the GEMM operations (GEMM1, corresponding to $\mathbf{P} \mathbf{V}$ for one iteration, and GEMM0, corresponding to $\mathbf{Q} \mathbf{K}^\top$ for the subsequent iteration) of warpgroup 1 precede those of warpgroup 2. Consequently, while warpgroup 2 is engaged in its GEMM operations, warpgroup 1 computes the softmax, and the two warpgroups subsequently alternate these roles—an approach referred to as ``pingpong'' scheduling, as visualized in~\cref{fig:pingpong_scheduling}. Although the scheduling pattern in practice exhibits more irregularity than the schematic suggests, we consistently observe notable performance benefits (for instance, FP16 forward pass throughput rises from 570 TFLOPS to between 620 and 640 TFLOPS for a head dimension of 128 and sequence length of 8192).

\paragraph{Attention variants} In handling multi-query attention~\citep{shazeer2019fast} and grouped query attention~\citep{ainslie2023gqa}, we adopt the methodology utilized in \textsc{FlashAttention-2}, modifying the tensor indexing scheme so that $\mathbf{K}$ and $\mathbf{V}$ are not redundantly stored in HBM.

\subsection{Intra-warpgroup overlapping GEMMs and softmax}

It is possible to achieve some instruction-level overlap between the softmax operations and the GEMMs within a single warpgroup. Below, we outline a method for accomplishing this.

Within the attention mechanism, the sequence of operations inside the main (inner) loop is characterized by dependencies between computations, restricting the possibility of parallel execution within a single loop iteration. As an illustration, the (local) softmax computation (from lines \ref{code-ws:softmax_start} to \ref{code-ws:softmax_end}) depends on the output $\mathbf{S}_i^{(j)}$ produced by the preceding GEMM operation, and subsequently, the next GEMM step requires as input the result $\widetilde{\mathbf{P}}_i^{(j)}$ computed by the softmax. This chain of dependencies is enforced by explicit wait statements at lines \ref{alg:ws_only_gemm1} and \ref{alg:ws_only_gemm2} in \cref{alg:flash3_wgmma_ws_only}, which enforce sequential execution of the GEMMs and the softmax routine. Nevertheless, by introducing pipelining across multiple loop iterations and incorporating supplementary buffers within registers, these serial dependencies can be decoupled. Based on this strategy, we introduce a two-stage\footnote{It is important to mention that the number of pipeline stages in the overlapping design is upper-bounded by, though does not have to match, the number $s$ of stages used in the circular SMEM buffer.} algorithm for pipelining the GEMM and softmax operations:

\begin{algorithm}[H]
    \caption{\small\label{alg:flash3_wgmma}\textsc{FlashAttention-3} consumer warpgroup forward pass}
    \begin{algorithmic}[1]
      \REQUIRE  Input matrices $\mathbf{Q}_i \in \mathbb{R}^{B_r \times d}$ as well as $\mathbf{K}, \mathbf{V} \in \mathbb{R}^{N \times d}$ located in HBM, and a specified key block size $B_c$ such that $T_c = \lceil \frac{N}{B_c} \rceil$.
      \STATE Dynamically repartition a fixed quantity of registers based on the current number of consumer warps.
      \STATE Allocate on-chip storage for $\mathbf{O}_i = (0) \in \mathbb{R}^{B_r \times d}$ as well as initialize $\ell_i, m_i = (0), (-\infty) \in \mathbb{R}^{B_r}$.
      \STATE Stall computation until both $\mathbf{Q}_i$ and $\mathbf{K}_0$ are available in shared memory.
      \STATE Using WGMMA, perform the computation $\mathbf{S}_{\mathrm{cur}} = \mathbf{Q}_i \mathbf{K}_0^T$. Commit the result and synchronize.
      \STATE Deallocate the $0$th stage buffer slot for $\mathbf{K}$.
      \STATE Derive $m_{i}$, $\tilde{\mathbf{P}}_{\mathrm{cur}}$, and $\ell_{i}$ from $\mathbf{S}_{\mathrm{cur}}$, then appropriately rescale $\mathbf{O}_i$.
      \FOR{$1 \le j < T_c - 1$}
        \STATE Ensure $\mathbf{K}_j$ has been fetched into shared memory.
        \label{code:mainloop_start}
        \STATE With WGMMA, calculate $\mathbf{S}_{\mathrm{next}} = \mathbf{Q}_i \mathbf{K}_{j}^T$, commit immediately without waiting. \label{code:first_wgmma}
        \STATE Confirm that $\mathbf{V}_{j-1}$ is present in shared memory.
        \STATE Employ WGMMA to accumulate $\mathbf{O}_{i} = \mathbf{O}_{i} + \tilde{\mathbf{P}}_{\mathrm{cur}} \mathbf{V}_{j-1}$; commit and do not wait. \label{code:second_wgmma}
        \STATE Block further operations until the result of WGMMA $\mathbf{Q}_i \mathbf{K}_{j}^T$ becomes available.
        \STATE Update $m_{i}$, compute $\tilde{\mathbf{P}}_{\mathrm{next}}$, and modify $\ell_{i}$ using $\mathbf{S}_{\mathrm{next}}$. \label{code:softmax}
        \STATE After ensuring $\tilde{\mathbf{P}}_{\mathrm{cur}} \mathbf{V}_{j-1}$ via WGMMA has finished, perform the necessary rescaling of $\mathbf{O}_i$
        \STATE Free the buffer slots: $(j\,\%\,s)$th for $\mathbf{K}$ and $(j-1\,\%\,s)$th for $\mathbf{V}$ accordingly.
        \STATE Propagate $\mathbf{S}_{\mathrm{next}}$ to $\mathbf{S}_{\mathrm{cur}}$.
        \label{code:mainloop_end}
      \ENDFOR \STATE Await the transfer of $\mathbf{V}_{T_c - 1}$ into shared memory.
      \STATE Using WGMMA, finalize computation with
      $\mathbf{O}_{i} = \mathbf{O}_{i} + \tilde{\mathbf{P}}_{\mathrm{last}} \mathbf{V}_{T_c - 1}$, commit and block until completed.

\STATE Epilogue: Adjust $\mathbf{O}_{i}$ according to $m_{i}$. Derive $L_{i}$ utilizing both $m_{i}$ and $\ell_{i}$.  
Store $\mathbf{O}_{i}$ and $L_{i}$ in HBM, placing them in the $i$-th position of $\mathbf{O}$ and $L$, respectively.

\cref{alg:flash3_wgmma} substitutes for the consumer pathway in \cref{alg:flash3_wgmma_ws_only} to yield the full \textsc{FlashAttention-3} procedure when using FP16. In essence, we employ the term WGMMA to refer to asynchronous GEMM. Throughout the principal loop (spanning lines \ref{code:mainloop_start} to \ref{code:mainloop_end}), each iteration $j$ overlaps the execution of the second WGMMA call (line \ref{code:second_wgmma}) with the softmax computation of the subsequent iteration $j+1$ (line \ref{code:softmax}).

Although the pipelined approach depicted previously suggests significant performance improvements in theory, several real-world considerations must be addressed:
\paragraph{Compiler reordering} The presented pseudocode assumes a specific execution order, yet in practice, the NVCC compiler may reorder instructions during optimization passes. This rearrangement can undermine the intended pipeline of alternating WGMMA and non-WGMMA instructions, potentially resulting in unanticipated behavior or reduced performance. However, inspecting the generated SASS code reveals that the compiler does produce the anticipated overlapping of instructions (as discussed in Section~\ref{sec:2-stage-sass}).
\paragraph{Register pressure} Achieving high efficiency requires minimizing register spills. The implementation of a 2-stage pipeline inherently raises register consumption, as it must hold extra intermediate states to facilitate the pipeline. In particular, it is necessary to allocate a register-resident buffer for $\mathbf{S}_{\mathrm{next}}$, thereby incurring an overhead of $B_r \times B_c \times \text{sizeof}(\text{float})$ additional register storage per threadblock. Such heightened register pressure may be at odds with large tile (block) sizes—which themselves demand substantial registers—forcing developers to weigh the benefits based on empirical profiling.
\paragraph{3-stage pipelining} Building upon the previous 2-stage methodology, we suggest a 3-stage pipeline that additionally overlaps the second WGMMA with the softmax operation. This modification promises to further improve Tensor Core usage. However, the inclusion of an extra pipeline stage exacerbates register requirements, making the decision between maximizing tile size and increasing pipeline depth even more challenging. Detailed exposition and evaluation of the 3-stage strategy is provided in ~\cref{sec:3-stage}.

\subsection{Low-precision with FP8}
\label{sec:algofp8}

\textbf{Efficiency: layout transformations.} Executing the forward computation for \textsc{FlashAttention-3} with FP8 precision introduces further complications regarding layout compatibility, which are not present when using FP16.

To begin with, it is important to recognize that the input tensors $\mathbf{Q}$, $\mathbf{K}$, and $\mathbf{V}$ are conventionally stored with contiguous memory layout along the head dimension. However, to comply with the k-major layout requirement imposed by FP8 WGMMA on the second GEMM, the tensor $\mathbf{V}$—specifically, the sub-tiles of $\mathbf{V}$ that are staged in SMEM—must be stored contiguously along the sequence length dimension. Because TMA load operations are unable to alter the dimension that is contiguous, this necessitates either (1) performing a transpose of $\mathbf{V}$ in GMEM before subsequent computation, or (2) conducting an in-kernel transpose of the relevant $\mathbf{V}$ tiles after they have been transferred to SMEM. The first strategy (1) can be realized either by (1a) integrating the transpose operation into the epilogue of a prior step, such as rotary embedding, or by (1b) launching a dedicated pre-processing transpose kernel\footnote{A well-tuned transpose kernel can achieve performance that approaches device bandwidth~\citep{colfax_cutlass_transpose_2024}.}, which interchanges the stride order between the sequence length and head dimensions. Nonetheless, approach (1a) presents significant integration challenges for generic library implementations, and (1b) is inefficient in memory-bound scenarios like inference due to additional overhead.

For FP8 \textsc{FlashAttention-3}, we select approach (2). Within the kernel, the transpose operation utilizes LDSM (\verb|ldmatrix|) and STSM (\verb|stmatrix|) instructions, which enable a warp of threads to jointly load shared memory (SMEM) into register memory (RMEM) and store from RMEM to SMEM, operating on 128-byte blocks.\footnote{The PTX specification indicates that LDSM and STSM function on $8 \times 8$ matrices with 16-bit elements \cite[\S 9.7.13.4.15-16]{ptx}, but by packing two 8-bit elements together, these instructions can be adapted for FP8 data. Nevertheless, the transpose variants of LDSM/STSM are unable to split packed 8-bit pairs, so explicit register-level data movement must occur between the LDSM and STSM phases to complete the matrix tile transpose; we leave these specifics aside.} These LDSM/STSM operations are efficient in their use of registers, permitting execution within the producer warpgroup, and inherently allow for layout transpositions during memory transfers. Furthermore, after the initial step, it is possible to schedule the transpose of the following $\mathbf{V}$ tile to overlap with the execution of the two WGMMAs operating on the prior $\mathbf{V}$ tile and the current $\mathbf{K}$ tile.

Second, we note a key distinction from FP16: the FP32 accumulator generated by an FP8 WGMMA possesses a memory layout that does not match the arrangement of operand A when it resides in registers. Illustrative excerpts of these respective layouts are shown in \cref{fig:rmem_accum} and \cref{fig:rmem_operand}, where elements are assigned to registers per thread in the order listed. By leveraging byte permutation instructions, it becomes feasible to reformat the output from the initial WGMMA accumulator to comply with the input expectations of a subsequent WGMMA, aligning with the structure of the $\mathbf{V}$ tile outputted by the in-kernel transpose. More concretely, as indicated in \cref{fig:rmem_accum}, the ordering is systematically altered to
$$\{ \verb|d0 d1 d4 d5 d2 d3 d6 d7| \},$$
and this pattern of register reordering is consistently repeated every 8 bytes. This manipulation effectively permutes the columns within the logical $\mathbf{P}$ tile; for instance, columns $0189$ are shifted to occupy the first four column slots. To ensure WGMMA produces the intended output tile, the in-kernel transpose may similarly be configured to emit the $\mathbf{V}$ tile with an analogous row reordering.\footnote{This enhanced flexibility due to performing the transpose within the kernel obviates the necessity of invoking shuffle instructions for transferring register ownership among threads, a strategy discussed previously in~\citep{colfax_fp8_flashattention_2024}.}

\textbf{Accuracy: block quantization and incoherent processing.} The FP8 (e4m3) format allocates just 3 bits for the mantissa and 4 bits to represent the exponent, which inherently leads to greater numerical inaccuracies compared to FP16/BF16 representations. Additionally, in large-scale models, it is common to encounter outlier values~\citep{dettmers2208llm, sun2024massive} whose magnitudes far exceed those of most other elements, presenting a significant challenge for quantization. A standard approach is per-tensor scaling~\citep{micikevicius2022fp8}, wherein a single scaling factor is maintained for each tensor (for instance, one each for $\mathbf{Q}$, $\mathbf{K}$, and $\mathbf{V}$). To mitigate the loss in accuracy during attention computations under FP8, we incorporate two specific strategies:

\begin{enumerate}

\item \textbf{Block quantization}: In this approach, a single scalar is assigned to represent each block; thus, for tensors $\mathbf{Q}$, $\mathbf{K}$, and $\mathbf{V}$, we partition them into blocks of dimensions either $B_r \times d$ or $B_c \times d$, performing quantization independently on each block. Notably, this quantization procedure can be merged with operations occurring immediately before the attention computation, such as rotary embeddings, without incurring any performance penalty (since these steps are constrained primarily by memory bandwidth). Furthermore, since the \textsc{FlashAttention-3} method natively works with blocks, we are able to rescale the blocks of $\mathbf{S}$ as necessary to adjust for the applied block quantization, all without extra computational overhead.

\item \textbf{Incoherent processing}: To mitigate the effect of outliers, we apply a random orthogonal transformation $\mathbf{M}$ to both $\mathbf{Q}$ and $\mathbf{K}$ before subjecting them to FP8 quantization. The orthogonality condition, $\mathbf{M} \mathbf{M}^\top = I$, ensures that $(\mathbf{Q} \mathbf{M})(\mathbf{K} \mathbf{M})^\top = \mathbf{Q} \mathbf{K}^\top$; as a result, this operation leaves the attention mechanism's output unaltered. The application of $\mathbf{M}$ helps to diffuse the influence of outlier values because each component of $\mathbf{Q} \mathbf{M}$ or $\mathbf{K} \mathbf{M}$ becomes a random combination of the original components, reducing the quantization error. Following the strategies of \citet{chee2024quip} and~\citet{tseng2024quip}, we construct $\mathbf{M}$ as the product of random diagonal matrices with entries $\pm1$ and a Hadamard matrix. This formulation allows for multiplication in $O(d \log d)$ time rather than $O(d^2)$, and like rotary embedding, it can be integrated into pre-attention computation without added computational expense.
\end{enumerate}

Our results, detailed in \cref{sec:numerical_error}, confirm that the adoption of these methods decreases numerical errors by as much as 2.6$\times$.

\section{Empirical Validation}
\label{sec:exp}

To implement \textsc{FlashAttention-3} and assess both its performance and precision, we employ CUTLASS~\citep{Thakkar_CUTLASS_2023} primitives, utilizing abstractions like WGMMA and TMA.

\begin{itemize}

\item \textbf{Benchmarking attention.} To assess performance, we evaluate the execution time of \textsc{FlashAttention-3} over varying sequence lengths and benchmark it against several baselines: the canonical PyTorch implementation, \textsc{FlashAttention-2}, a Triton version of \textsc{FlashAttention-2} (leveraging H100-targeted instructions), and a vendor-tuned \textsc{FlashAttention-2} variant built with cuDNN for H100 GPUs. Our results show that \textsc{FlashAttention-3} achieves up to a 2.0$\times$ speedup over \textsc{FlashAttention-2}, and is up to 1.5$\times$ faster than the Triton-based version. In terms of throughput, \textsc{FlashAttention-3} delivers up to 740 TFLOPs/s, corresponding to 75\% of the H100 GPU’s peak theoretical TFLOPs/s.
\item \textbf{Ablation study.} We demonstrate via ablation that both warp-specialization and the GEMM-softmax pipelining enhancements are key contributors to the observed acceleration of \textsc{FlashAttention-3}.
\item \textbf{Accuracy of FP8 attention.} Our analysis confirms that utilizing block quantization alongside incoherent processing leads to a 2.6$\times$ reduction in numerical errors for the FP8 variant of \textsc{FlashAttention-3}.
\end{itemize}

\subsection{Benchmarking Attention}
\label{subsec:benchmark_attn}

We evaluated the performance of various attention mechanisms on an H100 80GB SXM5 GPU across several configurations, including with and without a causal mask and using head dimensions of 64 or 128, processing inputs in FP16 format. The outcomes are presented in~\cref{fig:benchmark_attn_fwd} and~\cref{fig:benchmark_attn_bwd}. The data indicates that \textsc{FlashAttention-3} achieves forward pass speeds approximately 1.5-2.0$\times$ greater than those of \textsc{FlashAttention-2}, and backward pass speeds that are 1.5-1.75$\times$ higher. When compared with a conventional attention implementation, \textsc{FlashAttention-3} demonstrates improvements in speed of up to 3-16$\times$. Notably, for sequence lengths of 1,000 tokens or more, \textsc{FlashAttention-3} even outperforms the vendor-optimized cuDNN library, which is specifically tuned for H100 GPUs but remains closed source.

\paragraph{Benchmark settings:} The sequence length is adjusted among 512, 1k, ..., 16k, and the batch size is chosen such that the cumulative token count is 16k. The hidden dimension remains fixed at 2048, with the head dimension selected from 64, 128, or 256 (corresponding to 32, 16, or 8 attention heads, respectively). For computing the forward pass FLOPs, we use:

\begin{equation*}
  4 \cdot \text{seqlen}^2 \cdot \text{head dimension} \cdot \text{number of heads}.
\end{equation*}

When applying causal masking, this figure is halved to reflect that roughly 50\% of the elements are computed. The FLOPs for the backward pass are then obtained by multiplying the forward pass FLOPs by a factor of 2.5, since the backward pass involves 5 matrix multiplications compared to the 2 performed during the forward pass, as a result of recomputation.

Additionally, we evaluate the forward pass runtime using FP8 in comparable experimental conditions. The outcomes for headdim 256 are illustrated in ~\cref{fig:benchmark_attn_fp8}, while comprehensive findings are presented in \cref{sec:benchmark_attn_fp8_full}.

\subsection{Ablation Study: 2-Stage Pipelining Experiments}
\label{sec:2-stage-pipelining-experiments}

We conduct ablation studies on both the two-stage WGMMA-softmax pipelining and warp-specialization techniques for non-causal FP16 \textsc{FlashAttention-3}, using fixed parameters $\{\text{batch}, \text{seqlen}, \text{nheads}, \text{hdim}\} = \{ 4, 8448, 16, 128\}$. The findings presented in~\cref{table:ablation_pipelining} demonstrate that the introduced algorithmic enhancements—specifically, asynchrony through warp-specialization and concurrent execution of GEMM and softmax—yield a substantial increase in performance, raising throughput from 570 TFLOPs to 661 TFLOPs.

\subsection{Numerical Error Validation}
\label{sec:numerical_error}

Due to concerns raised about the numerical accuracy~\citep{golden2024flash} of \textsc{FlashAttention}, we evaluate and contrast the performance of \textsc{FlashAttention-2}, \textsc{FlashAttention-3}, and a standard attention algorithm, benchmarking each against an FP64 reference implementation. To replicate the presence of outlier activations and features observed in LLMs~\citep{dettmers2208llm, sun2024massive}, we sample the elements of $\mathbf{Q}, \mathbf{K}, \mathbf{V}$ from the distribution below:

\begin{equation*}
  \mathcal{N}(0, 1) + \mathcal{N}(0, 100) \cdot \mathrm{Bernoulli}(0.001).
\end{equation*}

Specifically, each matrix element follows a normal distribution with a mean of zero and a standard deviation of 1, except for 0.1\% of the elements, where an additional independent normal variable with a standard deviation of 10 is included. Subsequently, we compute the root mean squared error (RMSE), as shown in~\cref{table:numerical_error}. In FP16, both \textsc{FlashAttention-2} and \textsc{FlashAttention-3} report an RMSE that is 1.7$\times$ lower relative to the conventional implementation, owing to the use of FP32 for storing intermediate results (softmax). The baseline FP8 attention relies on per-tensor scaling, utilizes an FP32 accumulator for the matrix multiplication, and stores intermediate softmax computations in FP16. Owing to its use of block quantization and incoherent processing, \textsc{FlashAttention-3} in FP8 attains 2.6$\times$ higher accuracy than the baseline approach.

\section{Dicussion, Limitations, Conclusion}
\label{sec:discussion}

Through the development of \textsc{FlashAttention-3}, we have shown that leveraging innovative programming methods and harnessing hardware advancements—namely asynchrony and low-precision computation—can significantly enhance both the speed and precision of attention mechanisms. Our approach yields a 1.5-2.0$\times$ acceleration relative to \textsc{FlashAttention-2}, and achieves a 2.6$\times$ decrease in FP8 numerical error when compared with conventional per-tensor quantization methods. Nonetheless, there are several aspects for improvement that we intend to pursue, including better optimization for LLM inference, the incorporation of a persistent kernel design within the FP8 kernel,\footnote{In our reported results, only the FP16 \textsc{FlashAttention-3} implementation utilizes a persistent kernel and effective load balancing, whereas FP8 \textsc{FlashAttention-3} currently does not. This disparity partially accounts for the reduced performance of FP8 \textsc{FlashAttention-3} on small sequence lengths and when using causal masking, as opposed to the FP8 cuDNN kernels.} and a deeper investigation into how low-precision attention influences large-scale training scenarios. While our primary experiments are conducted on Hopper GPUs, we anticipate that these techniques will generalize well to a broad array of hardware accelerators. Ultimately, we believe that providing a more rapid and accurate attention primitive has the potential to enable new breakthroughs in long-context modeling tasks.

\end{document}